% Options for packages loaded elsewhere
\PassOptionsToPackage{unicode}{hyperref}
\PassOptionsToPackage{hyphens}{url}
%
\documentclass[
  10pt,
]{article}
\usepackage{amsmath,amssymb}
\usepackage{iftex}
\ifPDFTeX
  \usepackage[T1]{fontenc}
  \usepackage[utf8]{inputenc}
  \usepackage{textcomp} % provide euro and other symbols
\else % if luatex or xetex
  \usepackage{unicode-math} % this also loads fontspec
  \defaultfontfeatures{Scale=MatchLowercase}
  \defaultfontfeatures[\rmfamily]{Ligatures=TeX,Scale=1}
\fi
\usepackage{lmodern}
\ifPDFTeX\else
  % xetex/luatex font selection
\fi
% Use upquote if available, for straight quotes in verbatim environments
\IfFileExists{upquote.sty}{\usepackage{upquote}}{}
\IfFileExists{microtype.sty}{% use microtype if available
  \usepackage[]{microtype}
  \UseMicrotypeSet[protrusion]{basicmath} % disable protrusion for tt fonts
}{}
\makeatletter
\@ifundefined{KOMAClassName}{% if non-KOMA class
  \IfFileExists{parskip.sty}{%
    \usepackage{parskip}
  }{% else
    \setlength{\parindent}{0pt}
    \setlength{\parskip}{6pt plus 2pt minus 1pt}}
}{% if KOMA class
  \KOMAoptions{parskip=half}}
\makeatother
\usepackage{xcolor}
\usepackage[margin=0.75in]{geometry}
\usepackage{longtable,booktabs,array}
\usepackage{calc} % for calculating minipage widths
% Correct order of tables after \paragraph or \subparagraph
\usepackage{etoolbox}
\makeatletter
\patchcmd\longtable{\par}{\if@noskipsec\mbox{}\fi\par}{}{}
\makeatother
% Allow footnotes in longtable head/foot
\IfFileExists{footnotehyper.sty}{\usepackage{footnotehyper}}{\usepackage{footnote}}
\makesavenoteenv{longtable}
\setlength{\emergencystretch}{3em} % prevent overfull lines
\providecommand{\tightlist}{%
  \setlength{\itemsep}{0pt}\setlength{\parskip}{0pt}}
\setcounter{secnumdepth}{-\maxdimen} % remove section numbering
\ifLuaTeX
\usepackage[bidi=basic]{babel}
\else
\usepackage[bidi=default]{babel}
\fi
\babelprovide[main,import]{english}
% get rid of language-specific shorthands (see #6817):
\let\LanguageShortHands\languageshorthands
\def\languageshorthands#1{}
\ifLuaTeX
  \usepackage{selnolig}  % disable illegal ligatures
\fi
\IfFileExists{bookmark.sty}{\usepackage{bookmark}}{\usepackage{hyperref}}
\IfFileExists{xurl.sty}{\usepackage{xurl}}{} % add URL line breaks if available
\urlstyle{same}
\hypersetup{
  pdflang={en},
  hidelinks,
  pdfcreator={LaTeX via pandoc}}

\author{}
\date{}

\begin{document}

\hypertarget{the-correlated-anomaly-observation}{%
\section{The Correlated Anomaly
Observation}\label{the-correlated-anomaly-observation}}

Five independent measurement series, one fractional deviation band

\hypertarget{the-observation}{%
\subsection{1. The observation}\label{the-observation}}

Measurements made in four laboratories by three independent groups,
using four different detection techniques, over intervals from a single
solar flare to fifteen years, report fractional deviations that fall in
the same narrow band:

\textbackslash{[}10\^{}\{-4\}
\textbackslash;\textbackslash le\textbackslash;
\textbackslash frac\{\textbackslash delta\textbackslash lambda\}\{\textbackslash lambda\_0\}
\textbackslash;\textbackslash le\textbackslash;
10\^{}\{-3\}\textbackslash{]}

The series are unrelated to one another. They measure different
isotopes, at different facilities, with different instruments, for
different reasons. The shared quantity is the order of magnitude of the
deviation, not its sign, its phase, or its period.

\hypertarget{the-series}{%
\subsection{2. The series}\label{the-series}}

\begin{longtable}[]{@{}lllll@{}}
\toprule\noalign{}
Series & Facility & Interval & Fitted
\textbackslash(\textbackslash alpha\textbackslash) & Source \\
\midrule\noalign{}
\endhead
\bottomrule\noalign{}
\endlastfoot
\textbackslash(\^{}\{32\}\textbackslash mathrm\{Si\}\textbackslash)
seasonal & Brookhaven & 1982--1986, four years & \textbackslash(7.9
\textbackslash times 10\^{}\{-4\}\textbackslash) & Jenkins et al.
2008 \\
\textbackslash(\^{}\{36\}\textbackslash mathrm\{Cl\}\textbackslash)
seasonal & Brookhaven & 1982--1986, four years & \textbackslash(6.2
\textbackslash times 10\^{}\{-4\}\textbackslash) & Jenkins et al.
2008 \\
\textbackslash(\^{}\{226\}\textbackslash mathrm\{Ra\}\textbackslash)
seasonal & PTB, Germany & 1983--1998, fifteen years & \textbackslash(8.3
\textbackslash times 10\^{}\{-4\}\textbackslash) & Siegert et al.
1998 \\
\textbackslash(\^{}\{54\}\textbackslash mathrm\{Mn\}\textbackslash)
solar flare transient & Purdue & 13 December 2006, hours &
\textbackslash(\textbackslash alpha\textbackslash,A\_\{\textbackslash mathrm\{flare\}\}
= -2.5 \textbackslash times 10\^{}\{-3\}\textbackslash) & Jenkins \&
Fischbach 2009 \\
\textbackslash(\^{}\{32\}\textbackslash mathrm\{Si\}\textbackslash)/\textbackslash(\^{}\{36\}\textbackslash mathrm\{Cl\}\textbackslash)
amplitude & Brookhaven & 239 data points & \textbackslash(7.9(3)
\textbackslash times 10\^{}\{-4\}\textbackslash) & Jenkins et al.
2008 \\
\end{longtable}

Fitted coupling \textbackslash(\textbackslash alpha\textbackslash) under
the model of the time-gradient field, except where noted. The
\textbackslash(\^{}\{54\}\textbackslash mathrm\{Mn\}\textbackslash)
entry is not a bare \textbackslash(\textbackslash alpha\textbackslash):
the source reports the product
\textbackslash(\textbackslash alpha\textbackslash,A\_\{\textbackslash mathrm\{flare\}\}\textbackslash),
the coupling multiplied by the flare amplitude, and that product is
negative. It is listed as published and is not a like-for-like value in
the column. The Brookhaven series additionally reports an annual
oscillation of \textbackslash(7.9(3) \textbackslash times
10\^{}\{-4\}\textbackslash) with a 35(2) day phase shift, a correlation
of \textbackslash(r = 0.52\textbackslash) against
\textbackslash(1/R\^{}2\textbackslash), and a formal probability of
random correlation \textbackslash(p = 6 \textbackslash times
10\^{}\{-18\}\textbackslash). The
\textbackslash(\^{}\{54\}\textbackslash mathrm\{Mn\}\textbackslash)
transient preceded electromagnetic arrival by 36--40 hours.

\hypertarget{what-the-values-share}{%
\subsection{3. What the values share}\label{what-the-values-share}}

The three seasonal fitted couplings span \textbackslash(6.2
\textbackslash times 10\^{}\{-4\}\textbackslash) to \textbackslash(8.3
\textbackslash times 10\^{}\{-4\}\textbackslash), a ratio of
\textbackslash(1.34\textbackslash) across three isotopes and two
facilities. The table above lists four seasonal rows, but the fourth
repeats the Brookhaven amplitude already given in the first row rather
than adding a separate measurement, so there are three distinct seasonal
series. The
\textbackslash(\^{}\{54\}\textbackslash mathrm\{Mn\}\textbackslash)
transient
\textbackslash(\textbackslash alpha\textbackslash,A\_\{\textbackslash mathrm\{flare\}\}
= -2.5 \textbackslash times 10\^{}\{-3\}\textbackslash) is a factor of
about three larger in magnitude and is not a bare coupling. Across
different isotopes, different facilities, and different physical
circumstances, the couplings agree to within an order of magnitude.

No isotope-specific parameter was introduced to achieve this. The model
form is fixed in advance:

\textbackslash{[}\textbackslash lambda\_\{\textbackslash mathrm\{eff\}\}(t)
= \textbackslash lambda\_0 \textbackslash left( 1 + \textbackslash alpha
\textbackslash, G(t) \textbackslash right)\textbackslash{]}

with a single coupling
\textbackslash(\textbackslash alpha\textbackslash) per isotope and a
field \textbackslash(G(t)\textbackslash) of order
\textbackslash(10\^{}\{-3\}\textbackslash). The agreement across the
four series is therefore a constraint on the model form, not a
consequence of fitting freedom in it.

\hypertarget{reported-frequencies}{%
\subsection{4. Reported frequencies}\label{reported-frequencies}}

Power-spectrum analysis of the Brookhaven data by Sturrock et al. (2010)
reports the annual frequency, a prominent peak at
\textbackslash(11.18\textbackslash,\textbackslash text\{yr\}\^{}\{-1\}\textbackslash)
(a period of \textbackslash(32.7\textbackslash) days, conventionally
rounded to the \textbackslash(33\textbackslash)-day period) in the solar
synodic rotation band, and a running-mean peak at
\textbackslash(11.93\textbackslash,\textbackslash text\{yr\}\^{}\{-1\}\textbackslash).
The
\textbackslash(12.5\textbackslash,\textbackslash text\{yr\}\^{}\{-1\}\textbackslash)
frequency cited elsewhere in this body of work is reported by Sturrock,
Javorsek, Fischbach, Jenkins, and Lee (arXiv:1301.3754) as a periodicity
common to the BNL decay data and Super-Kamiokande solar neutrino data,
described there as indicative of the solar radiative zone. The r-mode
association in that literature belongs to the Rieger periodicity near
\textbackslash(2.1\textbackslash,\textbackslash text\{yr\}\^{}\{-1\}\textbackslash),
not to the
\textbackslash(12.5\textbackslash,\textbackslash text\{yr\}\^{}\{-1\}\textbackslash)
line.

\hypertarget{companion-series-in-other-domains}{%
\subsection{5. Companion series in other
domains}\label{companion-series-in-other-domains}}

The same band appears in two non-nuclear series, reported in the same
corpus:

\begin{longtable}[]{@{}lll@{}}
\toprule\noalign{}
Series & Fractional deviation & Basis \\
\midrule\noalign{}
\endhead
\bottomrule\noalign{}
\endlastfoot
Lunar Mascon residuals & \textbackslash(10\^{}\{-4\}\textbackslash) to
\textbackslash(10\^{}\{-3\}\textbackslash) & Residual after standard
spherical-harmonic modelling \\
Apollo navigation residuals & \textbackslash(10\^{}\{-4\}\textbackslash)
to \textbackslash(10\^{}\{-3\}\textbackslash) & Velocity and guidance
residuals on deep-space legs \\
\end{longtable}

\hypertarget{scope-of-this-page}{%
\subsection{6. Scope of this page}\label{scope-of-this-page}}

This page records an observation: that five independent measurement
series report deviations of the same fractional magnitude. It does not
claim a mechanism, and it does not present the field interpretation.

Every value above is quoted from published measurements or from the
companion papers in this body of work, with the source named in each
row. The interpretation of these deviations as a response to a scalar
field is developed separately, in the time-gradient field model and its
verification papers, and is not required to read this page.

\hypertarget{sources}{%
\subsection{7. Sources}\label{sources}}

{[}1{]} Jenkins, J.H. et al. "Evidence of Correlations between Nuclear
Decay Rates and the Earth-Sun Distance." \emph{Nuclear Physics A},
807(1--2), 81--101, 2008.

{[}2{]} Siegert, H., Schrader, H., and Schötzig, U. "Half-Life
Measurements of the Ra-226 Alpha Decay." \emph{Applied Radiation and
Isotopes}, 49(9--11), 1397--1401, 1998.

{[}3{]} Jenkins, J.H. and Fischbach, E. "Perturbation of nuclear decay
rates during the solar flare of 2006 December 13." \emph{Astroparticle
Physics}, 31(6), 407--411, 2009.

{[}4{]} Sturrock, P.A., Buncher, J.B., Fischbach, E., Gruenwald, J.T.,
Javorsek, D. II, Jenkins, J.H., Lee, R.H., Mattes, J.J., and Newport,
J.R. "Power Spectrum Analysis of BNL Decay-Rate Data."
\emph{Astroparticle Physics}, 34, 121--127, 2010.

{[}5{]} Sturrock, P.A., Javorsek, D., Fischbach, E., Jenkins, J.H., and
Lee, R.H. "The Case for a Solar Influence on Certain Nuclear Decay
Rates." \emph{arXiv:1301.3754}, 2013.

Full derivations, fitted parameters, and verification:
\href{https://richardkentgates.com/papers/time-gradient-field-model.html}{Time-Gradient
Field Model} and
\href{https://richardkentgates.com/papers/time_gradient_verification.html}{Mathematical
Verification of the Time-Gradient Field Model}.

\href{https://richardkentgates.com/}{← richardkentgates.com}

\hypertarget{ai-research-collaboration-disclosure}{%
\subsection{AI Research Collaboration
Disclosure}\label{ai-research-collaboration-disclosure}}

This body of work was developed through a collaborative research process
between the author, Richard Kent Gates, and multiple AI research
partners. The author provides full transparency on this process.

\textbf{Author\textquotesingle s Role (Richard Kent Gates):}

\begin{itemize}
\tightlist
\item
  Conceived all core theories, conceptual frameworks, hypotheses, and
  research direction
\item
  Defined the Time-Gradient Field Model, the unified scalar field G(t),
  the distance → time → information → G(t) framework, and all
  theoretical architecture
\item
  Provided physical intuition, biological reasoning, and domain
  expertise that guided every theoretical decision
\item
  Reviewed, validated, and approved all mathematical formalisms before
  inclusion
\item
  Made all final judgments on theoretical coherence and physical
  interpretation
\end{itemize}

\textbf{AI Research Partners:}

\begin{itemize}
\tightlist
\item
  \textbf{Mimo V2.5 (opencode platform)} --- Primary mathematical
  formalization partner. Assisted in translating conceptual physics into
  mathematical notation, generating numerical proof scripts, compiling
  LaTeX documents, and building the website infrastructure.
\item
  \textbf{Microsoft Copilot (browser-based)} --- Assisted in literature
  review, dataset gathering, cross-referencing empirical results (DESI,
  BNL, PTB, MICROSCOPE), and verifying citation accuracy.
\item
  \textbf{Google Gemini (browser-based)} --- Assisted in conceptual
  exploration, thought experimentation, and early-stage hypothesis
  refinement.
\item
  \textbf{Space Bunny (OpenCode)} --- Verification and provenance
  auditing. Read every paper and proof script in the repository and
  traced each reported figure back to the script and line that produces
  it. Independently recomputed the Fisher information matrices, the
  Cramér--Rao bounds, and the parameter errors from the models as
  stated. Resolved the origin of the Fisher information values quoted in
  the unified and Theory-of-Everything papers, reconstructed the
  two-parameter Fisher matrix in the statistical verification paper from
  its sampling grid, and identified a mislabeled citation and a
  statistical criterion stated more narrowly than the reported results
  supported. No theoretical claim, numerical result, or interpretive
  judgment in this body of work originated with this tool.
\end{itemize}

\textbf{Nature of AI Involvement:}

The AI tools functioned as research assistants --- analogous to graduate
students or technical collaborators who help formalize, compute, and
organize ideas that originate from the principal investigator. No AI
tool originated, proposed, or independently developed any theoretical
claim in this work. All physical insights, theoretical innovations, and
interpretive judgments are the author\textquotesingle s own.

\end{document}
